%% ECON 282E -- Session 7, Deck B3: Structure and Validation
%% Sources: AI-ECON Lec04_T4_Structured_NNs_Validation.tex, carried WHOLE
%% (27 frames); AI-ECON Lec03_T4 section "Specialized Architectures", carried
%% whole; Quant_Macro Lec.6 "Structured Neural Networks for Equilibrium
%% Models", "ICNN Example & Implementation" and "Monotonic Neural Networks";
%% Feng & Miao, "AI for Economists" (Book1), ch.10 "Embedding equilibrium
%% properties in the network" and "Accuracy validation and convergence
%% diagnostics".
%% Numbers from tools/figures/s07_numbers.json, keys "icnn", "five_methods",
%% "worked_nn_growth"; the classical rows are read from s05_numbers.json.

%% ECON 282E -- Foundations of Macroeconomics (UC Riverside, Fall 2026)
%% Shared Beamer preamble for every deck in slides/S01 ... slides/S10.
%%
%% Usage (a deck lives in slides/S0X/ and is compiled from that folder):
%%   %% ECON 282E -- Foundations of Macroeconomics (UC Riverside, Fall 2026)
%% Shared Beamer preamble for every deck in slides/S01 ... slides/S10.
%%
%% Usage (a deck lives in slides/S0X/ and is compiled from that folder):
%%   %% ECON 282E -- Foundations of Macroeconomics (UC Riverside, Fall 2026)
%% Shared Beamer preamble for every deck in slides/S01 ... slides/S10.
%%
%% Usage (a deck lives in slides/S0X/ and is compiled from that folder):
%%   \input{../preamble/course_preamble.tex}
%%   \decktitle{1}{A1}{Who I Am \& Why This Course}{September 24, 2026}
%%   \begin{document}
%%   \makedecktitle
%%   ...
%%
%% Adapted from Courses/AI-ECON-2026/source/shared_preamble.tex (Metropolis theme,
%% concept/caution/example/takeaway boxes, \sectiondivider). Changes for this course:
%%   * course title block (\decktitle / \makedecktitle) instead of \makebeamertitle
%%   * \zh{...} typesets Chinese (book titles) under pdflatex via CJKutf8
%%   * researchbox for the "Research in the wild" frames
%%   * section pages off; TikZ libraries and the course map loaded here
%% Build with pdflatex only (two passes); tools/build_slides.py does this for you.

\documentclass[10pt,english,aspectratio=169]{beamer}

\usepackage[T1]{fontenc}
\usepackage{textcomp}
\usepackage[utf8]{inputenc}
\usepackage{babel}
\usepackage{CJKutf8}

\usepackage{amsmath, amssymb, amsthm, graphicx}
\usepackage{booktabs, multirow, tabularx}
\usepackage{tikz}
\usetikzlibrary{positioning,arrows.meta,shapes,calc,fit,backgrounds}
\usepackage{pgfplots}
\pgfplotsset{compat=1.16}
\usepackage[most]{tcolorbox}
\tcbuselibrary{skins, breakable}
\usepackage{listings}
\usepackage{xcolor}

\lstset{
  basicstyle=\ttfamily\small,
  keywordstyle=\color{blue!80!black}\bfseries,
  commentstyle=\color{green!50!black}\itshape,
  stringstyle=\color{orange!80!black},
  backgroundcolor=\color{gray!8},
  frame=single,
  rulecolor=\color{gray!40},
  breaklines=true,
  showstringspaces=false,
  numbers=none,
}

\usetheme{metropolis}
\usefonttheme{professionalfonts}
\metroset{sectionpage=none}

\definecolor{LightGray}{RGB}{230,230,230}
\definecolor{RedTitle}{RGB}{0,0,200}
\definecolor{VeryLightGray}{RGB}{245,245,245}
\definecolor{DarkText}{RGB}{0,0,0}
\definecolor{UCRBlue}{RGB}{0,45,106}
\definecolor{UCRGold}{RGB}{241,171,0}

% Stage colours of the course arc (used by the course map and elsewhere)
\colorlet{StageTrunk}{blue!12}
\colorlet{StageBranch}{cyan!14}
\colorlet{StageMachine}{gray!18}
\colorlet{StageWall}{red!14}
\colorlet{StageTools}{orange!20}
\colorlet{StageData}{green!16}
\colorlet{StageAgents}{violet!14}

\hypersetup{bookmarks=false,breaklinks=true,pdfborder={0 0 0},colorlinks=true,
  linkcolor=DarkText,urlcolor=RedTitle}

\setbeamercolor{background canvas}{bg=white}
\setbeamercolor{title}{fg=RedTitle}
\setbeamercolor{frametitle}{bg=VeryLightGray,fg=DarkText}
\setbeamercolor{normal text}{fg=DarkText}
\setbeamercolor{alerted text}{fg=RedTitle}
\setbeamersize{text margin left=5mm, text margin right=5mm}
\addtobeamertemplate{frametitle}{}{\vspace{-0.5em}}

\setbeamertemplate{navigation symbols}{}
\setbeamertemplate{footline}{%
  \hspace{5mm}{\usebeamerfont{page number in head/foot}\color{black!45}%
  ECON 282E \,$\cdot$\, \insertshortdeck}\hfill%
  \usebeamercolor[fg]{page number in head/foot}%
  \usebeamerfont{page number in head/foot}%
  \insertframenumber\,/\,\inserttotalframenumber\kern1em\vskip2pt%
}

% ---------------------------------------------------------------------------
% Chinese text under pdflatex (book titles etc.)
\newcommand{\zh}[1]{\begin{CJK*}{UTF8}{gbsn}#1\end{CJK*}}

% ---------------------------------------------------------------------------
% Deck title block
%   \decktitle{session}{deck id}{title}{date}
\newcommand{\insertshortdeck}{}
\newcommand{\decksession}{}
\newcommand{\deckid}{}
\newcommand{\decktitletext}{}
\newcommand{\deckdate}{}
\newcommand{\decktitle}[4]{%
  \renewcommand{\decksession}{#1}%
  \renewcommand{\deckid}{#2}%
  \renewcommand{\decktitletext}{#3}%
  \renewcommand{\deckdate}{#4}%
  \renewcommand{\insertshortdeck}{Session #1 \,$\cdot$\, Deck #1.#2}%
  \title{#3}%
  \author{Zhigang Feng}%
  \date{#4}%
}
\newcommand{\makedecktitle}{%
  \begin{frame}[plain,noframenumbering]
    \begin{tikzpicture}[remember picture,overlay]
      \fill[UCRBlue] (current page.north west) rectangle ([yshift=-0.55cm]current page.north east);
      \fill[UCRGold] ([yshift=-0.55cm]current page.north west) rectangle ([yshift=-0.62cm]current page.north east);
      \node[anchor=west, text=white, font=\small\bfseries] at ([xshift=5mm,yshift=-0.275cm]current page.north west)
        {ECON 282E \,$\cdot$\, Foundations of Macroeconomics};
      \node[anchor=east, text=white, font=\small] at ([xshift=-5mm,yshift=-0.275cm]current page.north east)
        {UC Riverside \,$\cdot$\, Fall 2026};
    \end{tikzpicture}
    \vspace*{1.1cm}
    {\usebeamerfont{title}\color{black!55}\large Session \decksession\ \,$\cdot$\, Deck \decksession.\deckid\par}
    \vspace{0.25cm}
    {\usebeamerfont{title}\color{RedTitle}\LARGE\bfseries \decktitletext\par}
    \vspace{0.7cm}
    {\large Zhigang Feng\par}
    \vspace{0.15cm}
    {\color{black!65}\deckdate\par}
    \vfill
    {\footnotesize\itshape\color{black!55}Build it on paper $\;\to\;$ solve it on the machine $\;\to\;$ push it to the frontier with AI\par}
  \end{frame}}

% ---------------------------------------------------------------------------
% Section divider (plain frame, not counted)
\newcommand{\sectiondivider}[2]{%
\begin{frame}[plain,noframenumbering]
\begin{center}
\vfill
{\Large\textbf{#1}\par}
\vspace{0.35cm}
{\normalsize #2\par}
\vfill
\end{center}
\end{frame}
}

% ---------------------------------------------------------------------------
% Boxes
\newtcolorbox{conceptbox}[1][]{colback=blue!3, colframe=RedTitle!55, boxrule=0.35mm,
  arc=1mm, left=2mm, right=2mm, top=1mm, bottom=1mm, #1}
\newtcolorbox{cautionbox}[1][]{colback=red!3, colframe=red!50!black, boxrule=0.35mm,
  arc=1mm, left=2mm, right=2mm, top=1mm, bottom=1mm, #1}
\newtcolorbox{examplebox}[1][]{colback=green!3, colframe=green!45!black, boxrule=0.35mm,
  arc=1mm, left=2mm, right=2mm, top=1mm, bottom=1mm, #1}
\newtcolorbox{takeawaybox}[1][]{colback=VeryLightGray, colframe=RedTitle!55, boxrule=0.35mm,
  arc=1mm, left=2mm, right=2mm, top=1mm, bottom=1mm, #1}
\newtcolorbox{researchbox}[1][]{enhanced, colback=violet!4, colframe=violet!60!black,
  boxrule=0.35mm, arc=1mm, left=2mm, right=2mm, top=1mm, bottom=1mm,
  fonttitle=\bfseries\small, title={Research in the wild}, #1}

\newcommand{\compacttables}{\renewcommand{\arraystretch}{1.15}}
\newcommand{\src}[1]{{\tiny\color{black!55}#1}}

% ---------------------------------------------------------------------------
\input{../preamble/course_map.tex}

%%   \decktitle{1}{A1}{Who I Am \& Why This Course}{September 24, 2026}
%%   \begin{document}
%%   \makedecktitle
%%   ...
%%
%% Adapted from Courses/AI-ECON-2026/source/shared_preamble.tex (Metropolis theme,
%% concept/caution/example/takeaway boxes, \sectiondivider). Changes for this course:
%%   * course title block (\decktitle / \makedecktitle) instead of \makebeamertitle
%%   * \zh{...} typesets Chinese (book titles) under pdflatex via CJKutf8
%%   * researchbox for the "Research in the wild" frames
%%   * section pages off; TikZ libraries and the course map loaded here
%% Build with pdflatex only (two passes); tools/build_slides.py does this for you.

\documentclass[10pt,english,aspectratio=169]{beamer}

\usepackage[T1]{fontenc}
\usepackage{textcomp}
\usepackage[utf8]{inputenc}
\usepackage{babel}
\usepackage{CJKutf8}

\usepackage{amsmath, amssymb, amsthm, graphicx}
\usepackage{booktabs, multirow, tabularx}
\usepackage{tikz}
\usetikzlibrary{positioning,arrows.meta,shapes,calc,fit,backgrounds}
\usepackage{pgfplots}
\pgfplotsset{compat=1.16}
\usepackage[most]{tcolorbox}
\tcbuselibrary{skins, breakable}
\usepackage{listings}
\usepackage{xcolor}

\lstset{
  basicstyle=\ttfamily\small,
  keywordstyle=\color{blue!80!black}\bfseries,
  commentstyle=\color{green!50!black}\itshape,
  stringstyle=\color{orange!80!black},
  backgroundcolor=\color{gray!8},
  frame=single,
  rulecolor=\color{gray!40},
  breaklines=true,
  showstringspaces=false,
  numbers=none,
}

\usetheme{metropolis}
\usefonttheme{professionalfonts}
\metroset{sectionpage=none}

\definecolor{LightGray}{RGB}{230,230,230}
\definecolor{RedTitle}{RGB}{0,0,200}
\definecolor{VeryLightGray}{RGB}{245,245,245}
\definecolor{DarkText}{RGB}{0,0,0}
\definecolor{UCRBlue}{RGB}{0,45,106}
\definecolor{UCRGold}{RGB}{241,171,0}

% Stage colours of the course arc (used by the course map and elsewhere)
\colorlet{StageTrunk}{blue!12}
\colorlet{StageBranch}{cyan!14}
\colorlet{StageMachine}{gray!18}
\colorlet{StageWall}{red!14}
\colorlet{StageTools}{orange!20}
\colorlet{StageData}{green!16}
\colorlet{StageAgents}{violet!14}

\hypersetup{bookmarks=false,breaklinks=true,pdfborder={0 0 0},colorlinks=true,
  linkcolor=DarkText,urlcolor=RedTitle}

\setbeamercolor{background canvas}{bg=white}
\setbeamercolor{title}{fg=RedTitle}
\setbeamercolor{frametitle}{bg=VeryLightGray,fg=DarkText}
\setbeamercolor{normal text}{fg=DarkText}
\setbeamercolor{alerted text}{fg=RedTitle}
\setbeamersize{text margin left=5mm, text margin right=5mm}
\addtobeamertemplate{frametitle}{}{\vspace{-0.5em}}

\setbeamertemplate{navigation symbols}{}
\setbeamertemplate{footline}{%
  \hspace{5mm}{\usebeamerfont{page number in head/foot}\color{black!45}%
  ECON 282E \,$\cdot$\, \insertshortdeck}\hfill%
  \usebeamercolor[fg]{page number in head/foot}%
  \usebeamerfont{page number in head/foot}%
  \insertframenumber\,/\,\inserttotalframenumber\kern1em\vskip2pt%
}

% ---------------------------------------------------------------------------
% Chinese text under pdflatex (book titles etc.)
\newcommand{\zh}[1]{\begin{CJK*}{UTF8}{gbsn}#1\end{CJK*}}

% ---------------------------------------------------------------------------
% Deck title block
%   \decktitle{session}{deck id}{title}{date}
\newcommand{\insertshortdeck}{}
\newcommand{\decksession}{}
\newcommand{\deckid}{}
\newcommand{\decktitletext}{}
\newcommand{\deckdate}{}
\newcommand{\decktitle}[4]{%
  \renewcommand{\decksession}{#1}%
  \renewcommand{\deckid}{#2}%
  \renewcommand{\decktitletext}{#3}%
  \renewcommand{\deckdate}{#4}%
  \renewcommand{\insertshortdeck}{Session #1 \,$\cdot$\, Deck #1.#2}%
  \title{#3}%
  \author{Zhigang Feng}%
  \date{#4}%
}
\newcommand{\makedecktitle}{%
  \begin{frame}[plain,noframenumbering]
    \begin{tikzpicture}[remember picture,overlay]
      \fill[UCRBlue] (current page.north west) rectangle ([yshift=-0.55cm]current page.north east);
      \fill[UCRGold] ([yshift=-0.55cm]current page.north west) rectangle ([yshift=-0.62cm]current page.north east);
      \node[anchor=west, text=white, font=\small\bfseries] at ([xshift=5mm,yshift=-0.275cm]current page.north west)
        {ECON 282E \,$\cdot$\, Foundations of Macroeconomics};
      \node[anchor=east, text=white, font=\small] at ([xshift=-5mm,yshift=-0.275cm]current page.north east)
        {UC Riverside \,$\cdot$\, Fall 2026};
    \end{tikzpicture}
    \vspace*{1.1cm}
    {\usebeamerfont{title}\color{black!55}\large Session \decksession\ \,$\cdot$\, Deck \decksession.\deckid\par}
    \vspace{0.25cm}
    {\usebeamerfont{title}\color{RedTitle}\LARGE\bfseries \decktitletext\par}
    \vspace{0.7cm}
    {\large Zhigang Feng\par}
    \vspace{0.15cm}
    {\color{black!65}\deckdate\par}
    \vfill
    {\footnotesize\itshape\color{black!55}Build it on paper $\;\to\;$ solve it on the machine $\;\to\;$ push it to the frontier with AI\par}
  \end{frame}}

% ---------------------------------------------------------------------------
% Section divider (plain frame, not counted)
\newcommand{\sectiondivider}[2]{%
\begin{frame}[plain,noframenumbering]
\begin{center}
\vfill
{\Large\textbf{#1}\par}
\vspace{0.35cm}
{\normalsize #2\par}
\vfill
\end{center}
\end{frame}
}

% ---------------------------------------------------------------------------
% Boxes
\newtcolorbox{conceptbox}[1][]{colback=blue!3, colframe=RedTitle!55, boxrule=0.35mm,
  arc=1mm, left=2mm, right=2mm, top=1mm, bottom=1mm, #1}
\newtcolorbox{cautionbox}[1][]{colback=red!3, colframe=red!50!black, boxrule=0.35mm,
  arc=1mm, left=2mm, right=2mm, top=1mm, bottom=1mm, #1}
\newtcolorbox{examplebox}[1][]{colback=green!3, colframe=green!45!black, boxrule=0.35mm,
  arc=1mm, left=2mm, right=2mm, top=1mm, bottom=1mm, #1}
\newtcolorbox{takeawaybox}[1][]{colback=VeryLightGray, colframe=RedTitle!55, boxrule=0.35mm,
  arc=1mm, left=2mm, right=2mm, top=1mm, bottom=1mm, #1}
\newtcolorbox{researchbox}[1][]{enhanced, colback=violet!4, colframe=violet!60!black,
  boxrule=0.35mm, arc=1mm, left=2mm, right=2mm, top=1mm, bottom=1mm,
  fonttitle=\bfseries\small, title={Research in the wild}, #1}

\newcommand{\compacttables}{\renewcommand{\arraystretch}{1.15}}
\newcommand{\src}[1]{{\tiny\color{black!55}#1}}

% ---------------------------------------------------------------------------
%% Course map: the recurring "you are here" slide for ECON 282E.
%% Loaded by course_preamble.tex. Pattern follows AI-ECON-2026 lec01_map.tex, but the map
%% shows the whole ten-session course rather than one lecture.
%%
%%   \CourseMap{s}{label}   s = 1..10 highlights session s and prints "you are here: label"
%%                          (e.g. \CourseMap{1}{1.A2}); earlier sessions get a check mark.
%%   \CourseMap{0}{}        full overview, nothing highlighted.
%%
%% Note: TikZ option lists are comma-split before TeX conditionals run, so every \ifnum
%% below sits outside [...] option lists.

\newcommand{\CourseMap}[2]{%
\begin{center}
\begin{tikzpicture}[
    sess/.style={rectangle, rounded corners=2pt, draw=black!45, minimum width=1.34cm,
        minimum height=1.12cm, text width=1.26cm, align=center, inner sep=1pt},
    stage/.style={font=\tiny\bfseries, text=black!70},
]
  \foreach \i/\d/\t/\c in {%
      1/Sep 24/Intro \\ Growth/StageTrunk,
      2/Oct 1/HA $\cdot$ OLG \\ Policy/StageBranch,
      3/Oct 8/Num.\ DP \\ Python/StageMachine,
      4/Oct 15/Numerics \\ I/StageMachine,
      5/Oct 22/Numerics \\ II/StageMachine,
      6/Oct 29/The wall \\ \textbf{Midterm}/StageWall,
      7/Nov 5/ML \\ Deep nets/StageTools,
      8/Nov 12/RL \\ HA $+$ DL/StageTools,
      9/Nov 19/Text \\ LLMs/StageData,
      10/Dec 3/Agents \\ \textbf{Projects}/StageAgents} {
    \node[sess, fill=\c] (n\i) at ({1.46*(\i-1)}, 0)
      {{\scriptsize\bfseries S\i}\\[-1pt]{\tiny\color{black!60}\d}\\[1pt]{\tiny \t}};
    \ifnum\i<#1
      \node[font=\scriptsize, text=green!45!black] at ([xshift=-2.5mm,yshift=-2.2mm]n\i.north east) {$\checkmark$};
    \fi
    \ifnum\i=#1
      \draw[RedTitle, very thick, rounded corners=2pt]
        ([xshift=-0.6mm,yshift=-0.6mm]n\i.south west) rectangle ([xshift=0.6mm,yshift=0.6mm]n\i.north east);
      % keep the label inside the picture at both ends of the row
      \ifnum\i<3
        \node[font=\scriptsize\bfseries, text=RedTitle, anchor=south west] at ([yshift=1.2mm]n\i.north west)
          {$\blacktriangledown$\,you are here: #2};
      \else\ifnum\i>8
        \node[font=\scriptsize\bfseries, text=RedTitle, anchor=south east] at ([yshift=1.2mm]n\i.north east)
          {you are here: #2\,$\blacktriangledown$};
      \else
        \node[font=\scriptsize\bfseries, text=RedTitle, anchor=south] at ([yshift=1.2mm]n\i.north)
          {you are here: #2\,$\blacktriangledown$};
      \fi\fi
    \fi
  }
  % stage band under the sessions
  \foreach \a/\b/\lab/\c in {1/1/trunk/StageTrunk, 2/2/branches/StageBranch,
      3/5/the machine $\cdot$ classical methods/StageMachine, 6/6/the wall/StageWall,
      7/8/new tools/StageTools, 9/9/new data/StageData, 10/10/colleagues/StageAgents} {
    \fill[\c] ([yshift=-1.5mm]n\a.south west) rectangle ([yshift=-4.6mm]n\b.south east);
    \path ([yshift=-1.5mm]n\a.south west) -- ([yshift=-4.6mm]n\b.south east)
      node[midway, stage] {\lab};
  }
  \node[font=\footnotesize\itshape, text=black!70] at ([yshift=-9.5mm]$(n1.south)!0.5!(n10.south)$)
    {Build it on paper $\;\to\;$ solve it on the machine $\;\to\;$ push it to the frontier with AI};
\end{tikzpicture}
\end{center}
}


%%   \decktitle{1}{A1}{Who I Am \& Why This Course}{September 24, 2026}
%%   \begin{document}
%%   \makedecktitle
%%   ...
%%
%% Adapted from Courses/AI-ECON-2026/source/shared_preamble.tex (Metropolis theme,
%% concept/caution/example/takeaway boxes, \sectiondivider). Changes for this course:
%%   * course title block (\decktitle / \makedecktitle) instead of \makebeamertitle
%%   * \zh{...} typesets Chinese (book titles) under pdflatex via CJKutf8
%%   * researchbox for the "Research in the wild" frames
%%   * section pages off; TikZ libraries and the course map loaded here
%% Build with pdflatex only (two passes); tools/build_slides.py does this for you.

\documentclass[10pt,english,aspectratio=169]{beamer}

\usepackage[T1]{fontenc}
\usepackage{textcomp}
\usepackage[utf8]{inputenc}
\usepackage{babel}
\usepackage{CJKutf8}

\usepackage{amsmath, amssymb, amsthm, graphicx}
\usepackage{booktabs, multirow, tabularx}
\usepackage{tikz}
\usetikzlibrary{positioning,arrows.meta,shapes,calc,fit,backgrounds}
\usepackage{pgfplots}
\pgfplotsset{compat=1.16}
\usepackage[most]{tcolorbox}
\tcbuselibrary{skins, breakable}
\usepackage{listings}
\usepackage{xcolor}

\lstset{
  basicstyle=\ttfamily\small,
  keywordstyle=\color{blue!80!black}\bfseries,
  commentstyle=\color{green!50!black}\itshape,
  stringstyle=\color{orange!80!black},
  backgroundcolor=\color{gray!8},
  frame=single,
  rulecolor=\color{gray!40},
  breaklines=true,
  showstringspaces=false,
  numbers=none,
}

\usetheme{metropolis}
\usefonttheme{professionalfonts}
\metroset{sectionpage=none}

\definecolor{LightGray}{RGB}{230,230,230}
\definecolor{RedTitle}{RGB}{0,0,200}
\definecolor{VeryLightGray}{RGB}{245,245,245}
\definecolor{DarkText}{RGB}{0,0,0}
\definecolor{UCRBlue}{RGB}{0,45,106}
\definecolor{UCRGold}{RGB}{241,171,0}

% Stage colours of the course arc (used by the course map and elsewhere)
\colorlet{StageTrunk}{blue!12}
\colorlet{StageBranch}{cyan!14}
\colorlet{StageMachine}{gray!18}
\colorlet{StageWall}{red!14}
\colorlet{StageTools}{orange!20}
\colorlet{StageData}{green!16}
\colorlet{StageAgents}{violet!14}

\hypersetup{bookmarks=false,breaklinks=true,pdfborder={0 0 0},colorlinks=true,
  linkcolor=DarkText,urlcolor=RedTitle}

\setbeamercolor{background canvas}{bg=white}
\setbeamercolor{title}{fg=RedTitle}
\setbeamercolor{frametitle}{bg=VeryLightGray,fg=DarkText}
\setbeamercolor{normal text}{fg=DarkText}
\setbeamercolor{alerted text}{fg=RedTitle}
\setbeamersize{text margin left=5mm, text margin right=5mm}
\addtobeamertemplate{frametitle}{}{\vspace{-0.5em}}

\setbeamertemplate{navigation symbols}{}
\setbeamertemplate{footline}{%
  \hspace{5mm}{\usebeamerfont{page number in head/foot}\color{black!45}%
  ECON 282E \,$\cdot$\, \insertshortdeck}\hfill%
  \usebeamercolor[fg]{page number in head/foot}%
  \usebeamerfont{page number in head/foot}%
  \insertframenumber\,/\,\inserttotalframenumber\kern1em\vskip2pt%
}

% ---------------------------------------------------------------------------
% Chinese text under pdflatex (book titles etc.)
\newcommand{\zh}[1]{\begin{CJK*}{UTF8}{gbsn}#1\end{CJK*}}

% ---------------------------------------------------------------------------
% Deck title block
%   \decktitle{session}{deck id}{title}{date}
\newcommand{\insertshortdeck}{}
\newcommand{\decksession}{}
\newcommand{\deckid}{}
\newcommand{\decktitletext}{}
\newcommand{\deckdate}{}
\newcommand{\decktitle}[4]{%
  \renewcommand{\decksession}{#1}%
  \renewcommand{\deckid}{#2}%
  \renewcommand{\decktitletext}{#3}%
  \renewcommand{\deckdate}{#4}%
  \renewcommand{\insertshortdeck}{Session #1 \,$\cdot$\, Deck #1.#2}%
  \title{#3}%
  \author{Zhigang Feng}%
  \date{#4}%
}
\newcommand{\makedecktitle}{%
  \begin{frame}[plain,noframenumbering]
    \begin{tikzpicture}[remember picture,overlay]
      \fill[UCRBlue] (current page.north west) rectangle ([yshift=-0.55cm]current page.north east);
      \fill[UCRGold] ([yshift=-0.55cm]current page.north west) rectangle ([yshift=-0.62cm]current page.north east);
      \node[anchor=west, text=white, font=\small\bfseries] at ([xshift=5mm,yshift=-0.275cm]current page.north west)
        {ECON 282E \,$\cdot$\, Foundations of Macroeconomics};
      \node[anchor=east, text=white, font=\small] at ([xshift=-5mm,yshift=-0.275cm]current page.north east)
        {UC Riverside \,$\cdot$\, Fall 2026};
    \end{tikzpicture}
    \vspace*{1.1cm}
    {\usebeamerfont{title}\color{black!55}\large Session \decksession\ \,$\cdot$\, Deck \decksession.\deckid\par}
    \vspace{0.25cm}
    {\usebeamerfont{title}\color{RedTitle}\LARGE\bfseries \decktitletext\par}
    \vspace{0.7cm}
    {\large Zhigang Feng\par}
    \vspace{0.15cm}
    {\color{black!65}\deckdate\par}
    \vfill
    {\footnotesize\itshape\color{black!55}Build it on paper $\;\to\;$ solve it on the machine $\;\to\;$ push it to the frontier with AI\par}
  \end{frame}}

% ---------------------------------------------------------------------------
% Section divider (plain frame, not counted)
\newcommand{\sectiondivider}[2]{%
\begin{frame}[plain,noframenumbering]
\begin{center}
\vfill
{\Large\textbf{#1}\par}
\vspace{0.35cm}
{\normalsize #2\par}
\vfill
\end{center}
\end{frame}
}

% ---------------------------------------------------------------------------
% Boxes
\newtcolorbox{conceptbox}[1][]{colback=blue!3, colframe=RedTitle!55, boxrule=0.35mm,
  arc=1mm, left=2mm, right=2mm, top=1mm, bottom=1mm, #1}
\newtcolorbox{cautionbox}[1][]{colback=red!3, colframe=red!50!black, boxrule=0.35mm,
  arc=1mm, left=2mm, right=2mm, top=1mm, bottom=1mm, #1}
\newtcolorbox{examplebox}[1][]{colback=green!3, colframe=green!45!black, boxrule=0.35mm,
  arc=1mm, left=2mm, right=2mm, top=1mm, bottom=1mm, #1}
\newtcolorbox{takeawaybox}[1][]{colback=VeryLightGray, colframe=RedTitle!55, boxrule=0.35mm,
  arc=1mm, left=2mm, right=2mm, top=1mm, bottom=1mm, #1}
\newtcolorbox{researchbox}[1][]{enhanced, colback=violet!4, colframe=violet!60!black,
  boxrule=0.35mm, arc=1mm, left=2mm, right=2mm, top=1mm, bottom=1mm,
  fonttitle=\bfseries\small, title={Research in the wild}, #1}

\newcommand{\compacttables}{\renewcommand{\arraystretch}{1.15}}
\newcommand{\src}[1]{{\tiny\color{black!55}#1}}

% ---------------------------------------------------------------------------
%% Course map: the recurring "you are here" slide for ECON 282E.
%% Loaded by course_preamble.tex. Pattern follows AI-ECON-2026 lec01_map.tex, but the map
%% shows the whole ten-session course rather than one lecture.
%%
%%   \CourseMap{s}{label}   s = 1..10 highlights session s and prints "you are here: label"
%%                          (e.g. \CourseMap{1}{1.A2}); earlier sessions get a check mark.
%%   \CourseMap{0}{}        full overview, nothing highlighted.
%%
%% Note: TikZ option lists are comma-split before TeX conditionals run, so every \ifnum
%% below sits outside [...] option lists.

\newcommand{\CourseMap}[2]{%
\begin{center}
\begin{tikzpicture}[
    sess/.style={rectangle, rounded corners=2pt, draw=black!45, minimum width=1.34cm,
        minimum height=1.12cm, text width=1.26cm, align=center, inner sep=1pt},
    stage/.style={font=\tiny\bfseries, text=black!70},
]
  \foreach \i/\d/\t/\c in {%
      1/Sep 24/Intro \\ Growth/StageTrunk,
      2/Oct 1/HA $\cdot$ OLG \\ Policy/StageBranch,
      3/Oct 8/Num.\ DP \\ Python/StageMachine,
      4/Oct 15/Numerics \\ I/StageMachine,
      5/Oct 22/Numerics \\ II/StageMachine,
      6/Oct 29/The wall \\ \textbf{Midterm}/StageWall,
      7/Nov 5/ML \\ Deep nets/StageTools,
      8/Nov 12/RL \\ HA $+$ DL/StageTools,
      9/Nov 19/Text \\ LLMs/StageData,
      10/Dec 3/Agents \\ \textbf{Projects}/StageAgents} {
    \node[sess, fill=\c] (n\i) at ({1.46*(\i-1)}, 0)
      {{\scriptsize\bfseries S\i}\\[-1pt]{\tiny\color{black!60}\d}\\[1pt]{\tiny \t}};
    \ifnum\i<#1
      \node[font=\scriptsize, text=green!45!black] at ([xshift=-2.5mm,yshift=-2.2mm]n\i.north east) {$\checkmark$};
    \fi
    \ifnum\i=#1
      \draw[RedTitle, very thick, rounded corners=2pt]
        ([xshift=-0.6mm,yshift=-0.6mm]n\i.south west) rectangle ([xshift=0.6mm,yshift=0.6mm]n\i.north east);
      % keep the label inside the picture at both ends of the row
      \ifnum\i<3
        \node[font=\scriptsize\bfseries, text=RedTitle, anchor=south west] at ([yshift=1.2mm]n\i.north west)
          {$\blacktriangledown$\,you are here: #2};
      \else\ifnum\i>8
        \node[font=\scriptsize\bfseries, text=RedTitle, anchor=south east] at ([yshift=1.2mm]n\i.north east)
          {you are here: #2\,$\blacktriangledown$};
      \else
        \node[font=\scriptsize\bfseries, text=RedTitle, anchor=south] at ([yshift=1.2mm]n\i.north)
          {you are here: #2\,$\blacktriangledown$};
      \fi\fi
    \fi
  }
  % stage band under the sessions
  \foreach \a/\b/\lab/\c in {1/1/trunk/StageTrunk, 2/2/branches/StageBranch,
      3/5/the machine $\cdot$ classical methods/StageMachine, 6/6/the wall/StageWall,
      7/8/new tools/StageTools, 9/9/new data/StageData, 10/10/colleagues/StageAgents} {
    \fill[\c] ([yshift=-1.5mm]n\a.south west) rectangle ([yshift=-4.6mm]n\b.south east);
    \path ([yshift=-1.5mm]n\a.south west) -- ([yshift=-4.6mm]n\b.south east)
      node[midway, stage] {\lab};
  }
  \node[font=\footnotesize\itshape, text=black!70] at ([yshift=-9.5mm]$(n1.south)!0.5!(n10.south)$)
    {Build it on paper $\;\to\;$ solve it on the machine $\;\to\;$ push it to the frontier with AI};
\end{tikzpicture}
\end{center}
}


\graphicspath{{./figures/}}

%% Deck-local: the shared preamble is never edited (its md5 marks every deck stale).
\newcommand{\repo}[2]{\href{https://github.com/vonzhg/Macro-2026Fall/blob/main/#1}{\texttt{#2}}}

\lstset{language=Python, basicstyle=\ttfamily\scriptsize, frame=none,
        backgroundcolor=\color{gray!8}, aboveskip=2pt, belowskip=2pt}

\decktitle{7}{B3}{Structure, and How You Know It Worked}{Thursday, November 5, 2026}

\begin{document}

\makedecktitle

%=============================================================================
\begin{frame}{Where We Are}
\CourseMap{7}{7.B3}
\vspace{-1mm}
\begin{center}
\small \textbf{Block B:} 7.B1 the model as a loss function $\;\cdot\;$ 7.B2 deep VFI and actor--critic $\;\cdot\;$ \textbf{7.B3} structure and validation.
\end{center}
\end{frame}

%=============================================================================
\begin{frame}{Two Debts}
\begin{columns}[T]
\begin{column}{0.54\textwidth}
\small
Block B has produced two working solvers and two unpaid debts.
\medskip

\textbf{First, the networks obey no economics.} Theory says a value function is concave in capital and a consumption policy is increasing in wealth. Nothing in 7.B1 or 7.B2 imposes either. The network may satisfy them approximately where it was trained and violate them anywhere else --- and a violated shape restriction is not a small error, it is a \emph{qualitatively wrong} object.
\medskip

\textbf{Second, the \emph{a priori} bound is gone.} 3.A2 could bound the distance to the true value function \textbf{before looking at the answer}. Nothing in Block B can do that.
\end{column}
\begin{column}{0.43\textwidth}
\begin{conceptbox}[title=\textbf{This deck pays both}]\small
\textbf{Structure} --- build the restriction into the architecture, so it holds \emph{identically} rather than approximately. 7.A4's \textbf{fourth} constraint class.
\medskip
\textbf{Validation} --- replace the lost error bound with a measurement you can actually make.
\end{conceptbox}
\medskip
\begin{takeawaybox}\footnotesize
The course's first deck said it: \textbf{a numerical solution is a claim.} This is where that claim gets adjudicated for a network.
\end{takeawaybox}
\end{column}
\end{columns}
\end{frame}

%=============================================================================
\begin{frame}{Injecting Economic Structure}
\begin{columns}[T]
\begin{column}{0.54\textwidth}
\small
A generic feedforward network is \emph{deliberately} assumption-free: 7.A4 called a dense $\mathbf{W}$ the refusal to assume anything about the input. For images that refusal costs parameters. \textbf{In economics it costs something worse}, because we are not ignorant --- theory tells us the shape.
\medskip

\begin{itemize}\small\setlength{\itemsep}{1mm}
  \item $V(k)$ is \textbf{concave} in capital, from concave $u$ and $f$;
  \item $c(a)$ is \textbf{increasing} in wealth;
  \item a demand function is \textbf{monotone} in price;
  \item a cost function is \textbf{convex} in input prices.
\end{itemize}
\medskip

Every one of these is a restriction that can be built into the parameterisation instead of hoped for.
\end{column}
\begin{column}{0.43\textwidth}
\begin{takeawaybox}[title=\textbf{7.A4's fourth class}]\small
CNN: $\mathbf{W}$ sparse and repeated. RNN: $\mathbf{W}$ shared across dates. GNN: $\mathbf{W}$ follows the adjacency.
\medskip
\textbf{Here: constrain the \emph{signs} of the weights}, and the function is convex or monotone by construction.
\end{takeawaybox}
\medskip
\begin{examplebox}\footnotesize
And it is the same trade as ever. True, it buys parameters, data and reliability. False, it is bias --- which is why the restriction must come from theory, not from convenience.
\end{examplebox}
\end{column}
\end{columns}
\end{frame}

%=============================================================================
\begin{frame}{Measured: What an Unconstrained Network Does to Convexity}
\begin{columns}[T]
\begin{column}{0.52\textwidth}
\small
Fit a plain ReLU network ($1{,}153$ parameters) to a function that is \textbf{exactly convex} --- $0.6x^{2}+0.25x+0.1$, second derivative $1.2$ everywhere --- at \textbf{5 random seeds}.
\smallskip

\begin{center}\footnotesize
\begin{tabularx}{\linewidth}{@{}>{\raggedright\arraybackslash}X >{\raggedleft\arraybackslash}p{1.9cm} >{\raggedleft\arraybackslash}p{2.6cm}@{}}
\toprule
& \textbf{median RMSE} & \textbf{worst $-f''$} \\
\midrule
plain network & $\mathbf{0.00172}$ & $5.93$ \ \tiny(range $0$--$10.4$) \\
ICNN & $0.01345$ & $\mathbf{<7.1\times10^{-11}}$ \\
\bottomrule
\end{tabularx}
\end{center}
\smallskip
``Worst $-f''$'' is the \textbf{most negative curvature the fitted function reaches}: zero means it came out convex everywhere. At a typical seed the plain network fits the \emph{level} well and is \textbf{locally concave by about $5.9$}, against a true curvature of $1.2$.
\end{column}
\begin{column}{0.45\textwidth}
\begin{cautionbox}[title=\textbf{Fatal, not untidy}]\small
A concave patch breaks 7.B2's inner maximisation: local maxima appear, ascent lands in the wrong one, the policy \textbf{jumps}. \textbf{And it is invisible in the loss.}
\end{cautionbox}
\smallskip
\begin{takeawaybox}\footnotesize
\textbf{Read the range, not the median.} At one seed in five the plain network came out convex anyway: \textbf{not reliably broken, but unreliable} --- and you cannot tell which you have without checking. The ICNN never exceeds $7.1\times10^{-11}$.
\smallskip
\textbf{Not free:} the ICNN's RMSE is $\mathbf{7.8\times}$ larger. \emph{That is the price of the guarantee.}
\end{takeawaybox}
\end{column}
\end{columns}
\vspace{1mm}
\src{Ledger \texttt{icnn}, five seeds. ``Worst $-f''$'' is the largest negative second difference; zero means convex. Code: \repo{tools/figures/s07_generate.py}{s07\_generate.py}, \texttt{run\_icnn}; \repo{labs/L07_ML_and_Deep_Growth.ipynb}{L07} Part 6.}
\end{frame}

%=============================================================================
\begin{frame}{Input Convex Neural Networks}
\begin{columns}[T]
\begin{column}{0.54\textwidth}
\small
The architecture is a feedforward network with one extra channel and two restrictions:
\[
  z_{l+1}=\sigma_l\Bigl(W^{(z)}_{l}z_l+W^{(x)}_{l}x+b_l\Bigr).
\]
\medskip

\textbf{Sufficient for $f(x)$ to be convex in $x$:}
\begin{enumerate}\small\setlength{\itemsep}{0.7mm}
  \item the \textbf{recurrent} weights are non-negative, $W^{(z)}_l\ge0$;
  \item the activations $\sigma_l$ are \textbf{convex and non-decreasing} --- ReLU and softplus both are.
\end{enumerate}
\medskip

\textbf{The direct input weights $W^{(x)}_l$ are unconstrained}, including negative: every layer sees $x$ again through a free linear map, so the network can place and tilt the function freely.
\end{column}
\begin{column}{0.43\textwidth}
\begin{conceptbox}[title=\textbf{The proof, by induction}]\small
\textbf{Base.} $z_0=x$ is linear, hence convex.
\medskip
\textbf{Step.} $z_l$ convex; $W^{(z)}_l\ge0$ makes $W^{(z)}_lz_l$ a \textbf{non-negative combination of convex functions}. Adding the affine term preserves convexity, and so does applying a convex \emph{non-decreasing} $\sigma$.
\end{conceptbox}
\medskip
{\footnotesize Both conditions are load-bearing: a non-monotone activation such as $\tanh$ breaks the induction at the last step.}
\end{column}
\end{columns}
\vspace{1mm}
\src{Amos, Xu \& Kolter (2017), \emph{ICML}, ``Input Convex Neural Networks''.}
\end{frame}

%=============================================================================
\begin{frame}{Does the Constraint Cost Expressivity?}
\begin{columns}[T]
\begin{column}{0.53\textwidth}
\small
The obvious worry: with $W^{(z)}\ge0$, is universal approximation lost?
\medskip

\begin{conceptbox}[title=\textbf{Theorem (Chen, Shi \& Zhang, 2019)}]\small
An input convex neural network can approximate \textbf{any convex function} on a compact domain to arbitrary accuracy.
\end{conceptbox}
\medskip

The intuition is one line: a maximum of affine functions $\max_i\{a_i^{\top}x+b_i\}$ is convex, and \emph{any} convex shape is a limit of such maxima. An ICNN with ReLU units is a hierarchical max-affine approximator.
\end{column}
\begin{column}{0.44\textwidth}
\begin{takeawaybox}[title=\textbf{The reading that matters}]\small
You are \textbf{not} imposing a functional form --- not CES, not Cobb--Douglas, not translog.
\medskip
You are imposing a \textbf{shape restriction}, and keeping full non-parametric flexibility inside it.
\medskip
\emph{Architecture encodes the theory; the data choose the function within it.}
\end{takeawaybox}
\medskip
\begin{examplebox}\footnotesize
Exactly the situation an economist is usually in: theory pins the \emph{global shape} and says nothing about the curvature.
\end{examplebox}
\end{column}
\end{columns}
\vspace{1mm}
\src{Chen, Shi \& Zhang (2019), \emph{ICLR}, ``Optimal Control via Neural Networks: A Convex Approach''.}
\end{frame}

%=============================================================================
\begin{frame}{And What the Constraint Refuses}
\begin{columns}[T]
\begin{column}{0.53\textwidth}
\small
Point the same ICNN at $\sin 3x$ on $[-1,1]$, which is \textbf{not} convex, and train it identically.
\medskip

\begin{center}\footnotesize
\begin{tabularx}{\linewidth}{@{}>{\raggedright\arraybackslash}X >{\raggedleft\arraybackslash}p{2.2cm}@{}}
\toprule
\textbf{ICNN on $\sin 3x$} & \\
\midrule
RMSE, median of 5 & $0.3449$ \\
convexity violation & $<8.9\times10^{-11}$ \\
\bottomrule
\end{tabularx}
\end{center}
\medskip

It fails, completely --- and it fails while remaining \textbf{exactly convex}. The architecture did not bend; it could not.
\end{column}
\begin{column}{0.44\textwidth}
\begin{takeawaybox}[title=\textbf{This is the guarantee, from the other side}]\small
A constraint that can be violated when the data push hard enough is not a guarantee. \textbf{This one cannot be}, which is the entire point.
\medskip
The cost is stated exactly: if the true function is not convex, the ICNN will not find it, and the residual will say so loudly.
\end{takeawaybox}
\medskip
{\footnotesize And that is a diagnostic worth having: \textbf{a structured network that cannot fit is telling you the theory you imposed is wrong} --- an informative failure, unlike a silently wrong fit.}
\end{column}
\end{columns}
\vspace{1mm}
\src{Ledger \texttt{icnn.nonconvex}; same architecture, optimizer and budget as the convex case. Code: \repo{tools/figures/s07_generate.py}{s07\_generate.py}, \texttt{run\_icnn}.}
\end{frame}

%=============================================================================
\begin{frame}{What Convexity Buys a Macro Solver}
\begin{columns}[T]
\begin{column}{0.53\textwidth}
\small
\textbf{1. The inner maximisation becomes convex.} 7.B2's obstacle was $\max_{k'}\{u(k,k')+\beta\hat V(k')\}$ with $\hat V$ a non-convex network. With $\hat V$ concave by construction the objective is a \textbf{single-basin concave program}: gradient methods find the global maximum every time --- no bracketing, no multi-start, no policy flips.
\medskip

\textbf{2. The contraction comes back.} The Bellman operator preserves concavity, so an ICNN value function keeps $T$ a $\beta$-contraction on the relevant space. \textbf{Geometric convergence is restored}, at the classical rate.
\medskip

\textbf{3. Prices behave.} Convex preferences give monotone demand, hence a unique, stable market-clearing price in a calibration loop.
\end{column}
\begin{column}{0.44\textwidth}
\begin{takeawaybox}[title=\textbf{Read benefit 2 again}]\small
7.B2's comparison table had one genuinely alarming row: \emph{no general convergence guarantee}.
\medskip
\textbf{Imposing the shape the theory already gives you is what buys it back.} That is a far better reason to use an ICNN than accuracy.
\end{takeawaybox}
\medskip
\begin{examplebox}\footnotesize
\textbf{The cost is real and measured:} the constrained network fitted $7.8\times$ less accurately here. You are buying a \emph{guarantee} with \emph{fit} --- and for a Bellman solver that is a good trade.
\end{examplebox}
\end{column}
\end{columns}
\end{frame}

%=============================================================================
\begin{frame}{Partial Convexity, Because Macro Is Not That Tidy}
\begin{columns}[T]
\begin{column}{0.53\textwidth}
\small
A value function $V(k,z)$ is typically:
\begin{itemize}\small\setlength{\itemsep}{0.5mm}
  \item \textbf{concave in the endogenous states} $k$ --- capital, assets;
  \item \textbf{of no particular shape in the exogenous states} $z$ --- TFP, income risk.
\end{itemize}
\medskip

A plain ICNN would impose convexity on $z$ too. The \textbf{partially input convex} network splits the inputs:
\[
  z_{l+1}=\sigma\Bigl(\underbrace{W^{(z)}_l z_l}_{\ge0}+\underbrace{W^{(u)}_l(u\circ v)}_{\text{interaction}}+\underbrace{W^{(v)}_l v}_{\text{free}}+b_l\Bigr),
\]
with $k$ down the constrained path and $z$ down an ordinary one.
\end{column}
\begin{column}{0.44\textwidth}
\begin{takeawaybox}[title=\textbf{The general principle}]\small
Impose the restriction \textbf{exactly where theory gives you one}, and nowhere else. Over-imposing is bias by another route --- and unlike under-imposing, it does not show up in the residual.
\end{takeawaybox}
\smallskip
\begin{cautionbox}[title=\textbf{Softplus, not clamp}]\footnotesize
Enforce $W^{(z)}\ge0$ by \textbf{reparameterising}, $W=\mathrm{softplus}(\tilde W)$, never by clamping: a clamped weight at the boundary gets \textbf{zero gradient} and is dead. 7.A2's dying ReLU by another door.
\end{cautionbox}
\end{column}
\end{columns}
\end{frame}

%=============================================================================
\begin{frame}{Monotone Networks, and the Conflict They Create}
\begin{columns}[T]
\begin{column}{0.52\textwidth}
\small
The same device gives monotonicity. If every weight along a path from input $x_j$ to the output is non-negative and the activations are non-decreasing, then $\partial f/\partial x_j\ge0$ \textbf{identically}.
\medskip

\textbf{Partial monotonicity} is what is usually wanted: increasing in wealth, unrestricted in the shock. Constrain the signs only on the paths from the monotone inputs.
\medskip

\textbf{Cases where theory says exactly this:} consumption increasing in cash on hand, saving in wealth, demand decreasing in own price, the value function increasing in the state.
\end{column}
\begin{column}{0.46\textwidth}
\begin{cautionbox}[title=\textbf{The optimization conflict}]\small
Sign constraints \textbf{shrink the parameter set} and bend the loss surface, so a constrained network often trains more slowly and to a \emph{higher} training loss.
\medskip
\textbf{Not a defect:} the unconstrained network reaches its lower loss partly by using functions the theory forbids.
\end{cautionbox}
\medskip
{\footnotesize Two routes: \textbf{constrain} the architecture (exact, costs flexibility) or \textbf{penalise} violations in the loss (flexible, only approximate).}
\end{column}
\end{columns}
\end{frame}

%=============================================================================
\begin{frame}{The Instrument You Gave Up}
\begin{columns}[T]
\begin{column}{0.54\textwidth}
\small
3.A2 proved that value-function iteration is a contraction, and handed you a bound that needs no knowledge of the answer:
\[
  \lVert V_n-V^{*}\rVert_\infty\le\frac{\beta^{n}}{1-\beta}\lVert V_1-V_0\rVert_\infty .
\]
You could stop when it was small enough, and \textbf{know} how wrong you were.
\medskip

\textbf{No training loss replaces it.} That number is measured on the states the optimizer just drew, and 7.A3's whole regularisation discipline exists because it does not mean what it appears to mean.
\medskip

So the honest question is: \textbf{what \emph{can} be measured?} --- and the answer, next frame, is a bound that needs no training loss at all.
\end{column}
\begin{column}{0.43\textwidth}
\begin{cautionbox}[title=\textbf{Three things a small loss is consistent with}]\small
\textbf{1.} A policy that is excellent where you sampled and arbitrary elsewhere --- 7.B1 measured two decades lost at the band edge.
\medskip
\textbf{2.} A policy satisfying a \emph{necessary} condition on an explosive path.
\medskip
\textbf{3.} An actor and a critic that agree with each other and with nothing else.
\end{cautionbox}
\end{column}
\end{columns}
\vspace{1mm}
\src{Feng \& Miao, \emph{AI for Economists}, ch.~10, ``Accuracy validation and convergence diagnostics''.}
\end{frame}

%=============================================================================
\begin{frame}{The Off-Sample Euler Residual}
\begin{columns}[T]
\begin{column}{0.54\textwidth}
\small
The workhorse diagnostic, and the one every number in this course uses. After training, evaluate
\[
  \log_{10}\bigl|\mathcal{R}(k_i,z_i)\bigr|,\qquad i=1,\dots,M,
\]
on states \textbf{drawn independently of training} --- a fine grid, or a fresh sample, and deliberately including the edges.
\medskip

It is \textbf{unit-free}: divided by $u'(c)$, it reads as the relative consumption error a household makes by following this policy --- $-3$ is one part in a thousand.
\smallskip

\begin{center}\footnotesize
\begin{tabularx}{\linewidth}{@{}>{\raggedright\arraybackslash}p{2.6cm} >{\raggedright\arraybackslash}X@{}}
\toprule
$\log_{10}|\mathcal{R}|<-3$ & acceptable \\
$\log_{10}|\mathcal{R}|<-5$ & good \\
$\log_{10}|\mathcal{R}|<-7$ & excellent \\
\bottomrule
\end{tabularx}
\end{center}
\end{column}
\begin{column}{0.43\textwidth}
\begin{takeawaybox}[title=\textbf{Why it is the right instrument}]\footnotesize
Computed from the \textbf{model's own equilibrium conditions}, it needs no true solution to compare against --- exactly the situation you are in when no classical method can solve the model.
\medskip
\textbf{And it is a bound, not a proxy.} Santos (2000): the policy error is at most a constant times the maximum residual, the constant set by the curvature of $u$. 3.A1 states it; 5.B1 measured it.
\end{takeawaybox}
\medskip
{\footnotesize \textbf{Report the maximum, not the mean.} A mean hides a failure concentrated in part of the region: 7.B1's mean was $-3.43$ and its maximum $-1.61$.}
\end{column}
\end{columns}
\vspace{1mm}
\src{The benchmark scale is Judd (1998), \emph{Numerical Methods in Economics}.}
\end{frame}

%=============================================================================
\begin{frame}{Measured: This Session's Solutions on Judd's Scale}
\vspace{-1mm}
\begin{center}\footnotesize
\begin{tabularx}{\linewidth}{@{}>{\raggedright\arraybackslash}X >{\raggedleft\arraybackslash}p{2.1cm} >{\raggedleft\arraybackslash}p{2.1cm} >{\raggedright\arraybackslash}p{2.8cm}@{}}
\toprule
& \multicolumn{2}{c}{$\log_{10}\max|\mathcal{R}|$, off-sample} & \\
\textbf{Method} & \textbf{whole band} & \textbf{near $k^{*}$} & \textbf{verdict} \\
\midrule
Chebyshev collocation (5.B6) & $-6.40$ & $-6.82$ & good, nearly excellent \\
perturbation, 2nd order (5.A6) & $-4.50$ & $-6.14$ & acceptable; good near $k^{*}$ \\
global VFI (4.B5) & $-2.51$ & $-2.82$ & below acceptable \\
\textbf{neural network (7.B1)} & $\mathbf{-1.61}$ & $-3.25$ & \textbf{below acceptable on the band} \\
\quad \emph{interior $90\%$ only} & $-2.54$ & --- & still below $-3$ \\
\bottomrule
\end{tabularx}
\end{center}
\medskip
\begin{columns}[T]
\begin{column}{0.49\textwidth}
\begin{cautionbox}\footnotesize
\textbf{Say it plainly: our network does not reach Judd's ``acceptable'' across the band.} It reaches $-3.25$ near the steady state and $-1.61$ at the edge.
\end{cautionbox}
\end{column}
\begin{column}{0.49\textwidth}
\begin{takeawaybox}\footnotesize
\textbf{Which is what reporting the number is for.} The published figure for this kind of solver is $-3.5$ to $-5$; ours is worse, and a slide that quoted the literature instead of the run would have hidden that.
\end{takeawaybox}
\end{column}
\end{columns}
\vspace{1mm}
\src{Ledger \texttt{five\_methods}, \texttt{worked\_nn\_growth}; classical rows read from \texttt{s05\_numbers.json}. Code: \repo{tools/figures/s07_generate.py}{s07\_generate.py}, \repo{tools/figures/s05_generate.py}{s05\_generate.py}.}
\end{frame}

%=============================================================================
\begin{frame}{Why Run It Anyway}
\begin{columns}[T]
\begin{column}{0.52\textwidth}
\small
At $d=1$ the network spent thousands of parameters and minutes of compute to come \textbf{last} --- the \emph{expected} answer, not a disappointment.
\smallskip

\textbf{A network is not a more accurate approximator. It is a differently scaling one.} 7.A2 measured it with the parameter counts matched:
\smallskip

\begin{center}\footnotesize
\begin{tabularx}{\linewidth}{@{}>{\raggedright\arraybackslash}X >{\raggedleft\arraybackslash}p{1.9cm} >{\raggedleft\arraybackslash}p{1.9cm}@{}}
\toprule
$\max|\text{err}|$ & \textbf{Chebyshev} & \textbf{network} \\
 & $17$ coef. & $16$ par. \\
\midrule
smooth target & $\mathbf{2.5\times10^{-11}}$ & $2.5\times10^{-1}$ \\
kinked target & $1.8\times10^{-2}$ & $\mathbf{2.0\times10^{-4}}$ \\
\bottomrule
\end{tabularx}
\end{center}
\end{column}
\begin{column}{0.45\textwidth}
\begin{conceptbox}[title=\textbf{Where the crossover is}]\footnotesize
7.A2's rates say it plainly: a basis fixed in advance is capped at $O(m^{-2/d})$ --- the exponent carries $d$ --- while Barron's rate for a network does \textbf{not}.
\medskip
So the gap closes as $d$ grows, then reverses. \textbf{The crossover is not at $d=1$; it is out where a basis can no longer be written down.}
\end{conceptbox}
\smallskip
{\footnotesize And the same reversal appears at $d=1$ the moment the function is \textbf{kinked} --- which is what a borrowing limit does to a policy.}
\end{column}
\end{columns}
\vspace{1mm}
\begin{takeawaybox}\footnotesize
\textbf{So the rule is not ``networks are worse''.} It is: \textbf{if a classical method can solve your model, use it} --- and reach for the network when the state is a distribution, or the function has kinks you cannot locate in advance. S8 is entirely the first case.
\end{takeawaybox}
\end{frame}

%=============================================================================
\begin{frame}{The Benchmark Hierarchy}
\begin{columns}[T]
\begin{column}{0.54\textwidth}
\small
Validate against the strongest benchmark available, and be explicit about which one you had.
\medskip

\begin{enumerate}\small\setlength{\itemsep}{1.2mm}
  \item \textbf{A closed form.} The strongest, and rare --- Brock--Mirman with $\delta=1$. Tests \emph{correctness}: a method that cannot recover $k'=\alpha\beta k^{\alpha}$ has a bug.
  \item \textbf{A trusted classical solution.} 5.B6's Chebyshev solve, itself measured. Tests \emph{accuracy} on a model with no closed form.
  \item \textbf{The model's own conditions.} Off-sample Euler residuals. The only rung available at the frontier.
  \item \textbf{Self-consistency.} Different seeds, widths, samplers agreeing. The weakest, and not nothing.
\end{enumerate}
\end{column}
\begin{column}{0.43\textwidth}
\begin{takeawaybox}[title=\textbf{The discipline}]\small
\textbf{Climb down the ladder, never up.} Establish correctness on a model where rung 1 exists, then carry the \emph{same code} to a model where only rung 3 does.
\medskip
This is what 1.A1 meant by grading a neural solver gate by gate against a classical solution before the government is allowed into the model.
\end{takeawaybox}
\medskip
\begin{examplebox}\footnotesize
S8 does exactly this: a neural heterogeneous-agent solver graded against a validated Aiyagari reference before any new economics is added.
\end{examplebox}
\end{column}
\end{columns}
\vspace{1mm}
\src{Feng \& Miao, \emph{AI for Economists}, ch.~10's layered diagnostic system.}
\end{frame}

%=============================================================================
\begin{frame}{What Even a Good Residual Will Not Catch}
\begin{columns}[T]
\begin{column}{0.52\textwidth}
\small
\begin{itemize}\small\setlength{\itemsep}{1.2mm}
  \item \textbf{A necessary condition satisfied on the wrong path.} The Euler residual says nothing about transversality. A small residual is compatible with an explosive path.
  \item \textbf{Shape violations that do not move the residual.} A value function with a small concave patch can have a fine residual and a broken maximisation.
  \item \textbf{Regions never visited.} The residual is only measured where you measured it.
  \item \textbf{A wrong equilibrium.} With multiple equilibria, every one of them has a zero residual.
\end{itemize}
\end{column}
\begin{column}{0.46\textwidth}
\begin{takeawaybox}[title=\textbf{So check the economics too}]\small
Simulate and look at the \textbf{ergodic distribution}: does capital settle where theory says?
\medskip
Check \textbf{shape} directly --- the sign of the fitted second derivative on a grid, which is the number reported three frames back.
\medskip
Check \textbf{aggregates} against any moment you believe.
\end{takeawaybox}
\medskip
\begin{cautionbox}\footnotesize
And vary the seed. A result that moves with the random seed is not a result --- which is why width-stability across seeds is reported as evidence at all.
\end{cautionbox}
\end{column}
\end{columns}
\vspace{1mm}
\src{Feng \& Miao, \emph{AI for Economists}, ch.~10, ``Special challenges in a deep-learning setting''.}
\end{frame}

%=============================================================================
\begin{frame}{One Model, Five Methods --- the Reprise}
\begin{columns}[T]
\begin{column}{0.53\textwidth}
\small
The arc closes where it opened. The same quarterly RBC model has now been solved by \textbf{grid VFI} (4.B5), \textbf{first-} and \textbf{second-order perturbation} (5.A6), \textbf{Chebyshev collocation} (5.B6), and \textbf{a neural network} (7.B1).
\medskip

On this model the ranking is unambiguous and unflattering to the newest method. \textbf{Every column scales differently}, and that is the only reason the last row exists.
\medskip

\textbf{The rule to leave with:} if a classical method can solve your model, use it. Reach for the network when the state is a distribution, or when $d$ is large enough that a basis cannot be written down.
\end{column}
\begin{column}{0.44\textwidth}
\begin{takeawaybox}[title=\textbf{What the network did buy, even here}]\small
\textbf{A mesh-free policy} --- 4.B5's grid quantisation is simply absent.
\medskip
\textbf{Feasibility by construction}, not by clipping.
\medskip
\textbf{A method whose cost in $d$ is polynomial}, which is the only one of the five of which that is true.
\end{takeawaybox}
\medskip
\begin{cautionbox}\footnotesize
And what it cost is narrower than it looks: the \textbf{a priori} guarantee, 3.A2's $\beta^{n}/(1-\beta)$. What survives is the \textbf{a posteriori} bound --- Santos (2000) --- and, with shape constraints, the contraction itself.
\end{cautionbox}
\end{column}
\end{columns}
\vspace{1mm}
\src{Ledger \texttt{five\_methods}. The comparison was the point of running the network on a problem it was always going to lose. Code: \repo{tools/figures/s07_generate.py}{s07\_generate.py}; \repo{labs/L07_ML_and_Deep_Growth.ipynb}{L07}.}
\end{frame}

%=============================================================================
\begin{frame}{Takeaways, and Where Session 8 Goes}
\begin{columns}[T]
\begin{column}{0.55\textwidth}
\begin{takeawaybox}\small
\begin{itemize}\small\setlength{\itemsep}{1.1mm}
  \item Theory gives shapes, and a generic network obeys none of them. \textbf{Measured over 5 seeds: a plain net fitted to an exactly convex function came out locally concave by $5.9$ at the median} --- and convex at one. Not reliably wrong: \textbf{unreliable}.
  \item An \textbf{ICNN} is convex at every seed and approximates \emph{any} convex function --- and fitted $\mathbf{7.8\times}$ less accurately. \textbf{The guarantee has a price.}
  \item The deepest reason to impose it is not accuracy: \textbf{concavity restores the Bellman contraction} that 7.B2 gave up.
  \item Validation is the \textbf{off-sample} maximum Euler residual, on a known scale.
\end{itemize}
\end{takeawaybox}
\end{column}
\begin{column}{0.42\textwidth}
\begin{conceptbox}[title=\textbf{Bridge to Session 8}]\small
Every model in this session had a \textbf{small} state: $k$, or $(k,z)$. The networks were never needed --- that was the point of measuring.
\medskip
\textbf{S8 removes the training wheels.} The state becomes a \emph{distribution}; the transition $H(\Gamma)$ is no longer known; and the methods of this session become the only ones available.
\medskip
7.A4 already named the bridge: \textbf{permutation invariance}.
\end{conceptbox}
\end{column}
\end{columns}
\end{frame}

\end{document}
